\documentclass[a4paper,12pt]{article}
\usepackage[spanish]{babel}
\usepackage[utf8]{inputenc}
\usepackage[T1]{fontenc}
\usepackage{graphicx} % Para insertar imágenes
\usepackage{amsmath, amssymb} % Símbolos matemáticos
\usepackage{hyperref} % Enlaces y referencias clicables
\usepackage{lmodern} % Fuente moderna
\usepackage{geometry} % Márgenes
\usepackage[most]{tcolorbox} % Cajitas para ejemplos y procedimientos
\geometry{margin=2.5cm}

% ==== Definición de estilos de caja ====
\tcbset{
  colback=white,
  colframe=black!60,
  sharp corners,
  boxrule=0.8pt,
  left=6pt,right=6pt,top=6pt,bottom=6pt
}

\newtcolorbox{procedimiento}[1][]{
  colback=green!5!white,
  colframe=green!60!black,
  title=\textbf{Procedimiento},
  fonttitle=\bfseries,
  #1
}

\newtcolorbox{ejemplo}[1][]{
  colback=blue!5!white,
  colframe=blue!60!black,
  title=\textbf{Ejemplo},
  fonttitle=\bfseries,
  #1
}

% ==== Documento principal ====
\title{Álgebra Lineal y Estructuras Matemáticas (ALEM) Resumen Tema 1}
\author{Alonso Hernández Robles (Grupo E)}
\date{20/10/2025}

\begin{document}

\maketitle

\tableofcontents
\newpage

% =============================
% SECCIÓN: CAMBIO DE BASE
% =============================

\section{Cambio de Base}

\subsection{Pasar un número de base decimal a base \(b \geq 2\)}

\begin{procedimiento}
Dado un número decimal \( N \) y una base \( b \geq 2 \), para expresar \( N \) en base \( b \), se siguen los siguientes pasos:

\begin{enumerate}
    \item Dividir \( N \) entre \( b \) y anotar el \textbf{resto}.
    \item Tomar el \textbf{cociente} y volver a dividirlo entre \( b \), anotando el nuevo resto.
    \item Repetir el proceso hasta que el cociente sea \( 0 \).
    \item El número en base \( b \) se obtiene escribiendo los restos en orden \textbf{inverso} (del último al primero).
\end{enumerate}
\end{procedimiento}

\begin{ejemplo}[title={Pasar \(46_{10}\) a base \(5\)}]
\[
\begin{array}{rcl}
46 \div 5 &=& 9 \text{ con resto } 1\\
9 \div 5 &=& 1 \text{ con resto } 4\\
1 \div 5 &=& 0 \text{ con resto } 1
\end{array}
\]

Por tanto:
\[
46_{10} = (141)_5
\]
\end{ejemplo}

\begin{ejemplo}[title={Pasar \(46_{10}\) a base \(2\)}]
\[
\begin{array}{rcl}
46 \div 2 &=& 23 \text{ con resto } 0\\
23 \div 2 &=& 11 \text{ con resto } 1\\
11 \div 2 &=& 5 \text{ con resto } 1\\
5 \div 2 &=& 2 \text{ con resto } 1\\
2 \div 2 &=& 1 \text{ con resto } 0\\
1 \div 2 &=& 0 \text{ con resto } 1
\end{array}
\]

Por tanto:
\[
46_{10} = (101110)_2
\]
\end{ejemplo}

\begin{ejemplo}[title={Pasar \(46_{10}\) a base \(16\)}]
\[
46 \div 16 = 2 \text{ con resto } 14 = E
\]

Por tanto:
\[
46_{10} = (2E)_{16}
\]
\end{ejemplo}

\subsection{Pasar un número de base \(b\ge2\) a base decimal}

\begin{procedimiento}
Dado un número \( (a_n a_{n-1} \ldots a_1 a_0)_b \), su valor en base decimal es:

\[
N_{10} = a_n b^n + a_{n-1} b^{n-1} + \cdots + a_1 b^1 + a_0 b^0
\]

\begin{enumerate}
    \item Identificar los dígitos y su posición (empezando desde la derecha con exponente 0).
    \item Multiplicar cada dígito por la potencia de la base que le corresponde.
    \item Sumar todos los resultados.
\end{enumerate}
\end{procedimiento}

\begin{ejemplo}[title={Pasar \((101110)_2\) a base decimal}]
\[
(101110)_2 = 1\cdot 2^5 + 0\cdot 2^4 + 1\cdot 2^3 + 1\cdot 2^2 + 1\cdot 2^1 + 0\cdot 2^0
\]
\[
= 32 + 0 + 8 + 4 + 2 + 0 = 46
\]
\[
(101110)_2 = 46_{10}
\]
\end{ejemplo}

\subsection{Ecuaciones con Bases}

\begin{procedimiento}[title=Encontrar la base]
En una ecuación con base desconocida, el objetivo es determinar el valor de la base \( b \) que hace cierta la igualdad.

\begin{enumerate}
    \item Identificar el \textbf{mayor dígito} que aparece en la ecuación.  
    Entonces, debe cumplirse que:
    \[
    b \geq (\text{mayor dígito}) + 1
    \]
    ya que ningún dígito puede ser igual o mayor que la base.

    \item Expresar todos los números en \textbf{base decimal} usando potencias de \( b \).  
    Por ejemplo:
    \[
    (231)_b = 2b^2 + 3b + 1
    \]

    \item Sustituir en la ecuación original y simplificar para obtener una ecuación polinómica en \( b \).

    \item Resolver la ecuación y comprobar que el valor obtenido cumple el criterio del paso 1.
\end{enumerate}
\end{procedimiento}

\begin{ejemplo}[title={Ejemplo 1: Encontrar la base}]
Resolver en qué base \( b \) se cumple la ecuación:
\[
(3)_b \cdot (4)_b = (22)_b
\]

El mayor dígito que aparece es 4, por tanto:
\[
b \geq 4 + 1 = 5
\]

Pasamos todo a base decimal:
\[
3 \cdot 4 = 2b + 2
\]
\[
12 = 2b + 2 \Rightarrow 2b = 10 \Rightarrow b = 5 \ge 5
\]

\[
\boxed{b = 5}
\]
\end{ejemplo}

\begin{ejemplo}[title={Ejemplo 2: Encontrar la base}]
Determinar la base \( b \) para la que:
\[
(231)_b = (170)_{b+1}
\]

El mayor dígito en la ecuación es 7, por tanto:
\[
b + 1 \geq 8 \Rightarrow b \geq 7
\]

Expresamos ambos lados en base decimal:
\[
2b^2 + 3b + 1 = (b+1)^2 + 7(b+1)
\]

Desarrollamos el lado derecho:
\[
2b^2 + 3b + 1 = (b^2 + 2b + 1) + (7b + 7)
\]
\[
2b^2 + 3b + 1 = b^2 + 9b + 8
\]

Simplificamos:
\[
b^2 - 6b - 7 = 0
\]

Resolvemos la ecuación cuadrática:
\[
b = \frac{6 \pm \sqrt{(-6)^2 - 4(1)(-7)}}{2} = \frac{6 \pm \sqrt{36 + 28}}{2} = \frac{6 \pm 8}{2}
\]

\[
b_1 = 7, \quad b_2 = -1
\]

Descartamos \( b = -1 \) porque:
\[
b \geq 2 \text{ (por definición de base) y además } b \geq 7
\]

Por tanto, la única solución válida es:
\[
\boxed{b = 7}
\]
\end{ejemplo}


\begin{procedimiento}[title=Encontrar los números o dígitos]
Cuando la base es conocida y los \textbf{dígitos} son las incógnitas, el procedimiento es el siguiente:

\begin{enumerate}
    \item Expresar los números en base decimal como suma de potencias de la base.  
    Por ejemplo, si la base es 5:
    \[
    (abc)_5 = a\cdot5^2 + b\cdot5 + c
    \]
    \item Sustituir en la ecuación y simplificar para obtener una ecuación con incógnitas \( a, b, c, \ldots \)
    \item Tener en cuenta que los dígitos siempre deben cumplir:
    \[
    a, b, c, \ldots \in \{0,1,2,\ldots,b-1\}
    \]
    \item Probar los valores posibles hasta encontrar los que satisfacen la ecuación.
\end{enumerate}
\end{procedimiento}

\begin{ejemplo}[title={Ejemplo: Encontrar número con cifras intercambiadas en base 16}]
Encuentre todos los números naturales de dos cifras en base 10 que, en base 16, se escriben con las mismas cifras pero en orden invertido.

\[
t = (ab)_{10} = (ba)_{16}, \quad a,b \in \{0,\dots,9\}
\]

Pasamos a base decimal para plantear la ecuación:
\[
10a + b = 16b + a
\]

\[
10a + b = 16b + a \Rightarrow 9a - 15b = 0 \Rightarrow 3a = 5b
\]

Como \(3a = 5b\), analizamos divisibilidad:

\begin{itemize}
    \item \(5 \mid 3a \Rightarrow 5 \mid a \Rightarrow a = 0,5\)
    \item \(3 \mid 5b \Rightarrow 3 \mid b \Rightarrow b = 0,3,6,9\)
\end{itemize}

Comprobamos los casos:

1. \(a = 0 \Rightarrow 3a = 0 \Rightarrow 5b = 0 \Rightarrow b = 0\)  
   \(\Rightarrow t_1 = 0\)

2. \(a = 5 \Rightarrow 3a = 15 \Rightarrow 5b = 15 \Rightarrow b = 3\)  
   \(\Rightarrow t_2 = 53\)

\medskip
\textbf{Conclusión:} Los números que cumplen la condición son:
\[
\boxed{t = 0 \text{ y } t = 53}
\]
\end{ejemplo}


% =============================
% RESTO DE SECCIONES
% =============================

\section{Propiedades de la Divisibilidad}

La divisibilidad es una relación fundamental entre números enteros.  
Decimos que un número entero \( a \) \textbf{divide} a otro número entero \( b \) si existe un entero \( k \) tal que \( b = ak \).  
En tal caso se escribe \( a \mid b \).

A continuación, se presentan algunas propiedades básicas de la divisibilidad en \( \mathbb{Z} \).

\begin{procedimiento}[title={Propiedad 1}]
Para todo \( a \in \mathbb{Z} \), se cumple que:
\[
1 \mid a \quad \text{y} \quad a \mid 0
\]
\end{procedimiento}

\begin{procedimiento}[title={Propiedad 2}]
Si \( a \mid b \) y \( b \mid a \), entonces:
\[
a = \pm b
\]
\end{procedimiento}

\begin{procedimiento}[title={Propiedad 3}]
Si \( a \mid b \), entonces para todo \( c \in \mathbb{Z} \):
\[
a \mid (b \cdot c)
\]
\end{procedimiento}

\begin{procedimiento}[title={Propiedad 4}]
Si \( a \mid b \) y \( b \mid c \), entonces:
\[
a \mid c
\]
\end{procedimiento}

\begin{procedimiento}[title={Propiedad 5}]
Si \( a \mid b \) y \( a \mid b' \), entonces:
\[
a \mid (b + b')
\]
\end{procedimiento}


% =============================
% SECCIÓN: MCD Y MCM
% =============================

\section{Propiedades del MCD y MCM}

El \textbf{máximo común divisor (MCD)} y el \textbf{mínimo común múltiplo (MCM)} son dos funciones aritméticas fundamentales que relacionan números enteros mediante la divisibilidad.

A continuación, se presentan algunas propiedades importantes de ambos.

% ====================================
% SUBSECCIÓN: MÁXIMO COMÚN DIVISOR
% ====================================

\subsection{Propiedades del MCD}

\begin{procedimiento}[title={Propiedad 1}]
\[
\text{mcd}(0,0,\ldots,0) = 0
\]
\end{procedimiento}

\begin{procedimiento}[title={Propiedad 2}]
El MCD de varios números puede calcularse de forma recursiva:
\[
\text{mcd}(a_1,a_2,\ldots,a_{n-1},a_n) = \text{mcd}(\text{mcd}(a_1,a_2,\ldots,a_{n-1}),a_n)
\]
\end{procedimiento}

\begin{procedimiento}[title={Propiedad 3}]
\[
\text{mcd}(a,b) = \text{mcd}(-a,b) = \text{mcd}(a,-b) = \text{mcd}(-a,-b)
\]
\end{procedimiento}

\begin{procedimiento}[title={Propiedad 4}]
\[
\text{mcd}(a,0) = |a|
\]
\end{procedimiento}

\begin{procedimiento}[title={Propiedad 5}]
Si \( a = bq + d \), entonces:
\[
\text{mcd}(a,b) = \text{mcd}(b,d)
\]
\end{procedimiento}

\begin{ejemplo}[title={Ejemplo de uso: Valores de $\text{mcd}(n^2-n+1,n+2)$ en $\mathbb{Z}$}]
Dividimos \(n^2 - n + 1\) entre \(n+2\):

\[
n^2-n+1 / n+2 = n-3 \text{ resto } 7
\]

\[
n^2 - n + 1 = (n+2)(n-3) + 7
\]

Por la Propiedad 5:
\[
\text{mcd}(n^2 - n + 1, n+2) = \text{mcd}(n+2,7) \le 7
\]

Como 7 es primo, los posibles valores son:
\[
\boxed{1 \text{ o } 7}
\]
\end{ejemplo}


% ====================================
% SUBSECCIÓN: MÍNIMO COMÚN MÚLTIPLO
% ====================================

\subsection{Propiedades del MCM}

\begin{procedimiento}[title={Propiedad 1}]
\[
\text{mcm}(0,a_2,a_3,\ldots,a_n) = 0
\]
\end{procedimiento}

\begin{procedimiento}[title={Propiedad 2}]
El MCM de varios números también puede calcularse de forma recursiva:
\[
\text{mcm}(a_1,a_2,\ldots,a_{n-1},a_n) = \text{mcm}(\text{mcm}(a_1,a_2,\ldots,a_{n-1}),a_n)
\]
\end{procedimiento}

\begin{procedimiento}[title={Propiedad 3}]
\[
\text{mcm}(a,b) = \text{mcm}(-a,b) = \text{mcm}(a,-b) = \text{mcm}(-a,-b)
\]
\end{procedimiento}

% =============================
% SECCIÓN: NÚMEROS PRIMOS
% =============================

\section{Números Primos}

Los \textbf{números primos} son la base de la aritmética entera.  
Su estudio es esencial para comprender la factorización, la divisibilidad y el cálculo del MCD y MCM.

% ====================================
% DEFINICIÓN
% ====================================

\begin{procedimiento}[title={Definición de número primo}]
Un número entero \( p \) es \textbf{primo} si cumple que:
\[
p \geq 2 \quad \text{y sus únicos divisores son } \pm 1 \text{ y } \pm p
\]
\end{procedimiento}

\begin{procedimiento}[title={Propiedad fundamental de los números primos}]
Si \( p \) es primo y \( p \mid a \cdot b \), entonces:
\[
p \mid a \quad \text{o} \quad p \mid b
\]
\end{procedimiento}

% ====================================
% CRITERIO DE PRIMALIDAD
% ====================================

\begin{procedimiento}[title={Cómo comprobar si un número es primo}]
Un número \( n \geq 2 \) es \textbf{primo} si no existe ningún número primo \( p \) tal que:
\[
p \leq \sqrt{n} \quad \text{y} \quad p \mid n
\]
\end{procedimiento}

\begin{ejemplo}[title={Ejemplo: comprobar si \(47\) es primo}]
Calculamos:
\[
\sqrt{47} \approx 6.85
\]
Primos menores o iguales que \(6\): \(2, 3, 5\)
\[
47 \div 2 = 23.5 \quad (\text{no entero})
\]
\[
47 \div 3 = 15.66 \quad (\text{no entero})
\]
\[
47 \div 5 = 9.4 \quad (\text{no entero})
\]
Ningún primo menor o igual que \(\sqrt{47}\) divide a \(47\), por tanto:
\[
\boxed{47 \text{ es primo.}}
\]
\end{ejemplo}

% ====================================
% FACTORIZACIÓN Y RELACIÓN ENTRE MCD Y MCM (fusionado)
% ====================================

\begin{procedimiento}[title={Factorización prima y fórmulas de MCD y MCM}]
Todo número entero distinto de cero puede escribirse (de forma única salvo el signo) como:
\[
a = \pm \displaystyle\prod_{i=1}^{n} P_i^{c_i}, \qquad 
b = \pm \displaystyle\prod_{i=1}^{n} P_i^{d_i}
\]
donde \( P_i \) son primos distintos y \( c_i, d_i \geq 0 \).

A partir de esta factorización, se cumplen simultáneamente las siguientes fórmulas:
\[
\text{mcd}(a,b) = \displaystyle\prod_{i=1}^{n} P_i^{\min(c_i, d_i)}
\]
\[
\text{mcm}(a,b) = \displaystyle\prod_{i=1}^{n} P_i^{\max(c_i, d_i)}
\]
\[
\text{mcd}(a,b) \cdot \text{mcm}(a,b) = |a \cdot b|
\]

\end{procedimiento}

\section{Algoritmo Extendido de Euclides y Coeficientes de Bézout}

\begin{procedimiento}[title={Pasos del Algoritmo Extendido de Euclides}]
Para calcular el \(\text{mcd}(a,b)\) y sus coeficientes de Bézout:

\begin{enumerate}
    \item Realizar divisiones sucesivas:
    \[
    a = b \cdot q_1 + r_1, \quad
    b = r_1 \cdot q_2 + r_2, \quad \ldots
    \]
    hasta que el resto sea \(0\). El último resto no nulo ($r_n$) es el \(\text{mcd}(a,b)\).

    \item Construir la tabla iterativa con columnas:
    \[
    i,\quad r_i,\quad q_i,\quad u_i,\quad v_i
    \]
    donde \(u_i, v_i\) se definen por la recurrencia:
    \[
    u_i = u_{i-2} - u_{i-1}\,q_i, \qquad
    v_i = v_{i-2} - v_{i-1}\,q_i,
    \]
    con condiciones iniciales:
    \[
    u_0 = 0, \quad u_1 = 1, \quad v_0 = 1, \quad v_1 = -q_1.
    \]

    \item La tabla queda entonces:
    \[
    \begin{array}{c|c|c|c|c}
    i & r_i & q_i & u_i & v_i \\
    \hline
    0 & b & - & 0 & 1 \\
    1 & r_1 & q_1 & 1 & -q_1 \\
    2 & r_2 & q_2 & u_2 & v_2 \\
    \vdots & \vdots & \vdots & \vdots & \vdots \\
    n & \boxed{r_n} & q_n & u_n & v_n \\
    \end{array}
    \]

    \item \(u_n\) y \(v_n\) verifican la identidad de Bézout:
    \[
    \text{mcd}(a,b) = a \cdot u_n + b \cdot v_n = \boxed{r_n}
    \]
\end{enumerate}
\end{procedimiento}


% =========================================================
% EJEMPLO 1: MCD(114,51)
% =========================================================

\begin{ejemplo}[title={Cálculo de \(\text{mcd}(114,51)\) y coeficientes de Bézout}]
\textbf{Divisiones:}

\[
114 / 51 = 2 \;\text{resto}\; 12
\]
\[
51 / 12 = 4 \;\text{resto}\; \boxed{3}
\]
\[
12 / 3 = 4 \;\text{resto}\; 0
\]

El último resto no nulo es \(\boxed{3}\).

\bigskip

\textbf{Tabla del Algoritmo Extendido:}

\[
\begin{array}{c|c|c|c|c}
i & r_i & q_i & u_i & v_i \\
\hline
0 & 51 & \text{-} & 0 & 1 \\
1 & 12 & 2 & 1 & -2 \\
2 & \boxed{3} & 4 & -4 & 9 \\
\end{array}
\]

\bigskip

\textbf{Coeficientes de Bézout:}

\[
\boxed{3} = 114 \cdot (-4) + 51 \cdot (9).
\]

Por tanto, \(-4\) y \(9\) son coeficientes de Bézout para \(114\) y \(51\).
\end{ejemplo}

% =========================================================
% EJEMPLO 2: MCD(6,10,15)
% =========================================================

\begin{ejemplo}[title={Cálculo de \(\text{mcd}(6,10,15)\)}]
Usamos la propiedad:
\[
\text{mcd}(6,10,15) = \text{mcd}(\text{mcd}(6,10),15).
\]

\subsubsection*{Paso 1: \(\text{mcd}(6,10)\)}

\textbf{Divisiones:}

\[
10 / 6 = 1 \;\text{resto}\; 4
\]
\[
6 / 4 = 1 \;\text{resto}\; \boxed{2}
\]
\[
4 / 2 = 2 \;\text{resto}\; 0
\]

El último resto no nulo es \(\boxed{2}\).

\bigskip

\textbf{Tabla del Algoritmo Extendido:}

\[
\begin{array}{c|c|c|c|c}
i & r_i & q_i & u_i & v_i \\
\hline
0 & 6 & \text{-} & 0 & 1 \\
1 & 4 & 1 & 1 & -1 \\
2 & \boxed{2} & 1 & -1 & 2 \\
\end{array}
\]

\[
\boxed{2} = 10 \cdot (-1) + 6 \cdot (2).
\]

\vspace{1em}

\subsubsection*{Paso 2: \(\text{mcd}(2,15)\)}

\textbf{Divisiones:}

\[
15 / 2 = 7 \;\text{resto}\; \boxed{1}
\]
\[
2 / 1 = 2 \;\text{resto}\; 0
\]

El último resto no nulo es \(\boxed{1}\).

\bigskip

\textbf{Tabla del Algoritmo Extendido:}

\[
\begin{array}{c|c|c|c|c}
i & r_i & q_i & u_i & v_i \\
\hline
0 & 2 & \text{-} & 0 & 1 \\
1 & \boxed{1} & 7 & 1 & -7 \\
\end{array}
\]

\[
\boxed{1} = 15 \cdot (1) + 2 \cdot (-7).
\]
\end{ejemplo}

\begin{tcolorbox}[colback=blue!5!white, colframe=blue!60!black, title=\textbf{Cálculo de \(\text{mcd}(6,10,15)\)}]
\textbf{Combinación de resultados:}
\[
\text{mcd}(6,10) = \boxed{2} = 10 \cdot (-1) + 6 \cdot (2),
\]
\[
\text{mcd}(2,15) = \boxed{1} = 15 \cdot (1) + 2 \cdot (-7).
\]

Sustituyendo \(2 = 10\cdot(-1) + 6\cdot(2)\) en la ecuación anterior:

\[
1 = 15 \cdot (1) + \big(10\cdot(-1) + 6\cdot(2)\big)\cdot(-7)
\]
\[
1 = 15\cdot(1) + 10\cdot(7) + 6\cdot(-14).
\]

Por tanto:
\[
1,\; 7,\; -14
\]
son unos coeficientes de Bézout para \(15,\;10,\;6\).
\end{tcolorbox}

\section{Ecuaciones Diofánticas}

\subsection{Con 2 incógnitas}

\begin{procedimiento}
Consideremos la ecuación diofántica:
\[
a x + b y = c, \quad a,b \neq 0.
\]

\begin{enumerate}
    \item Calcular \(g = \text{mcd}(a,b)\). Si \(g \nmid c\), la ecuación no tiene solución.
    \item Si \(g \mid c\), dividir la ecuación entre \(g\) para simplificar:
    \[
    \frac{a}{g} x + \frac{b}{g} y = \frac{c}{g}.
    \]
    \item Hallar una solución particular \((x_0,y_0)\) usando los coeficientes de Bézout de \(a\) y \(b\):
    \[
    x_0 \cdot a + y_0 \cdot b = g.
    \]
    \item Las soluciones generales son:
    \[
    x = x_0 + \frac{b}{g} k, \quad y = y_0 - \frac{a}{g} k, \quad k \in \mathbb{Z}.
    \]
\end{enumerate}
\end{procedimiento}


\begin{ejemplo}[title={Ecuación sin solución}]
\[
6x + 15y = 7
\]

\textbf{Divisiones para el mcd:}

\[
15 / 6 = 2 \;\text{resto}\; \boxed{3}
\]
\[
6 / 3 = 2 \;\text{resto}\; 0
\]

\[
\text{mcd}(6,15) = \boxed{3}
\]

Como \(3 \nmid 7\), \textbf{no tiene solución.}
\end{ejemplo}

\begin{ejemplo}[title={Ecuación con solución}]
\[
21x + 15y = 33
\]

\textbf{Divisiones para el mcd:}

\[
21 / 15 = 1 \;\text{resto}\; 6
\]
\[
15 / 6 = 2 \;\text{resto}\; \boxed{3}
\]
\[
6 / 3 = 2 \;\text{resto}\; 0
\]

\[
\text{mcd}(21,15) = \boxed{3} \quad \text{y } 3 \mid 33 \Rightarrow \text{tiene solución.}
\]

\textbf{Tabla del Algoritmo Extendido de Euclides:}

\[
\begin{array}{c|c|c|c|c}
i & r_i & q_i & u_i & v_i \\
\hline
0 & 15 & - & 0 & 1 \\
1 & 6 & 1 & 1 & -1 \\
2 & \boxed{3} & 2 & -2 & 3 \\
\end{array}
\]

\[
\boxed{3} = 21 \cdot (-2) + 15 \cdot (3)
\]

\textbf{Solución particular:}  
Multiplicamos la combinación de Bézout por \(33/3 = 11\):

\[
x_0 = -2 \cdot 11 = -22, \quad y_0 = 3 \cdot 11 = 33
\]

\textbf{Solución general:}

\[
x = -22 + \frac{15}{3} k = -22 + 5 k
\]
\[
y = 33 - \frac{21}{3} k = 33 - 7 k
\]

\[
k \in \mathbb{Z}
\]

\textbf{Observación:} la ecuación es equivalente a
\[
7x + 5y = 11
\]
dividida entre el mcd \(3\).
\end{ejemplo}

\subsection{Con 3 incógnitas}

\begin{procedimiento}
Consideremos la ecuación diofántica:
\[
a x + b y + c z = d, \quad a,b,c \neq 0.
\]

\begin{enumerate}
    \item Comprobar si tiene solución calculando
    \[
    g = \text{mcd}(a,b,c) = \text{mcd}(\text{mcd}(a,b),c).
    \]
    Si \(g \nmid d\), la ecuación no tiene solución.
    
    \item Si \(g \mid d\), dividimos toda la ecuación entre \(g\) para simplificar.
    
    \item Hallar los mcd particulares de los pares de coeficientes:
    \[
    \text{mcd}(a,b), \quad \text{mcd}(a,c), \quad \text{mcd}(b,c).
    \]
    
    \item Dependiendo de estos valores, distinguimos dos casos:
    
    \subsubsection*{Primer caso: al menos un par tiene mcd = 1}
    \begin{itemize}
        \item Identificar un par de coeficientes cuyo mcd sea 1.
        \item Dejar los términos de ese par en la izquierda y pasar el resto de términos junto con el término independiente a la derecha.
        \item Establecer la variable de la derecha como \(k \in \mathbb{Z}\) (por ejemplo, \(x = k\)) y resolver la ecuación de 2 incógnitas resultante.
        \item La solución general quedará expresada en función de dos parámetros enteros \(k\) y \(k' \in \mathbb{Z}\).
    \end{itemize}

    \subsubsection*{Segundo caso: todos los pares tienen mcd \(\neq 1\)}
    \begin{itemize}
        \item Elegir un par de coeficientes, por ejemplo \(a\) y \(b\), y factorizar su mcd:
        \[
        a x + b y = g \cdot t.
        \]
        \item Sustituir esta descomposición en la ecuación original, obteniendo una ecuación de dos incógnitas: \(g \cdot t + c z = d\).
        \item Resolver esta ecuación de dos incógnitas como se hizo previamente.
        \item Luego sustituir el valor solución de \(t\) en la identidad \(t = 2x + 5y\) (según el ejemplo, se ve mucho más fácil) y resolver la ecuación resultante de dos incógnitas.
        \item La solución final se expresará en función de dos parámetros enteros \(k\) y \(k' \in \mathbb{Z}\).
    \end{itemize}
\end{enumerate}
\end{procedimiento}

%%%%%%%%%%%%%%%%%%%%%%%%%%%%%%%%%%%%%%%%%%%%%%%%%%%%%%%%%%%%
\begin{ejemplo}[title={Ecuación diofántica con 3 incógnitas (Primer caso)}]
Consideremos la ecuación:
\[
10x + 15y - 14z = 8
\]

\textbf{Paso 1: Comprobación de existencia de soluciones}

\[
\text{mcd}(10,15,-14) = \text{mcd}(\text{mcd}(10,15),14)
\]

\[
15 / 10 = 1 \;\text{resto}\; \boxed{5}
\]
\[
10 / 5 = 2 \;\text{resto}\; 0
\]
\[
\Rightarrow \text{mcd}(10,15) = 5
\]

\[
\text{mcd}(5,14): \quad 14 / 5 = 2 \;\text{resto}\; 4
\]
\[
5 / 4 = 1 \;\text{resto}\; \boxed{1}
\Rightarrow \text{mcd}(5,14) = 1
\]

Por tanto, 
\[
\text{mcd}(10,15,-14) = 1.
\]
Como \(1 \mid 8\), la ecuación \textbf{tiene solución}.

\bigskip

\textbf{Paso 2: Determinar el tipo de caso}

\[
\begin{aligned}
\text{mcd}(10,15) &= 5 \neq 1, \\
\text{mcd}(10,-14) &= \text{mcd}(10,14) = 2 \neq 1, \\
\text{mcd}(15,-14) &= \text{mcd}(15,14) = 1.
\end{aligned}
\]

Como \(\text{mcd}(15,-14)=1\), estamos ante el \textbf{primer caso}.

\bigskip

\textbf{Paso 3: Reducción a dos incógnitas}

Dejamos a la izquierda los términos cuyo mcd es 1:
\[
15y - 14z = 8 - 10x.
\]

Sea \(x = k, \; k \in \mathbb{Z}\). Sustituyendo:
\[
15y - 14z = 8 - 10k.
\]

\bigskip

\textbf{Paso 4: Resolución de la ecuación de dos incógnitas}

\textit{Divisiones para hallar el mcd(15,14):}
\[
15 / 14 = 1 \;\text{resto}\; \boxed{1}
\]
\[
\Rightarrow \text{mcd}(15,14) = 1.
\]

\end{ejemplo}

\begin{ejemplo}[title={Ecuación diofántica con 3 incógnitas (Primer caso)}]

\textbf{Tabla del Algoritmo Extendido:}
\[
\begin{array}{c|c|c|c|c}
i & r_i & q_i & u_i & v_i \\
\hline
0 & 14 & \text{-} & 0 & 1 \\
1 & \boxed{1} & 1 & 1 & -1 \\
\end{array}
\]

\[
\boxed{1} = 15 \cdot (1) + 14 \cdot (-1).
\]

\bigskip

Reorganizamos y multiplicamos por \((8 - 10k)\):
\[
15 \cdot (8 - 10k) - 14 \cdot (8 - 10k) = 8 - 10k.
\]

Por tanto, una \textbf{solución particular} es:
\[
y_0 = 8 - 10k, \quad z_0 = 8 - 10k.
\]

\bigskip

\textbf{Paso 5: Solución general}

\[
\begin{aligned}
y &= y_0 + \frac{-14}{1} \, k' = 8 - 10k - 14k', \\
z &= z_0 - \frac{15}{1} \, k' = 8 - 10k - 15k',
\end{aligned}
\quad k' \in \mathbb{Z}.
\]

\bigskip

\textbf{Por tanto, las soluciones generales son:}
\[
\boxed{
\begin{aligned}
x &= k,\\
y &= 8 - 10k - 14k',\\
z &= 8 - 10k - 15k',
\end{aligned}
\quad k, k' \in \mathbb{Z}.
}
\]
\end{ejemplo}

%%%%%%%%%%%%%%%%%%%%%%%%%%%%%%%%%%%%%%%%%%%%%%%%%%%%%%%%%%%%
\begin{ejemplo}[title={Ecuación diofántica con 3 incógnitas (Segundo caso)}]
Consideremos la ecuación:
\[
6x + 15y + 10z = 7
\]

\textbf{Paso 1: Comprobación de existencia de soluciones}

\[
\text{mcd}(6,15,10) = \text{mcd}(\text{mcd}(6,15),10)
\]

\[
15 / 6 = 2 \;\text{resto}\; \boxed{3}
\]
\[
6 / 3 = 2 \;\text{resto}\; 0
\]
\[
\Rightarrow \text{mcd}(6,15) = 3
\]

\[
10 / 3 = 3 \;\text{resto}\; \boxed{1} \Rightarrow \text{mcd}(3,10) = 1
\]
Por tanto:
\[
\text{mcd}(6,15,10) = 1
\]
Como \(1 \mid 7\), la ecuación \textbf{tiene solución}.

\bigskip

\textbf{Paso 2: Determinar el tipo de caso}

\[
\begin{aligned}
\text{mcd}(6,15) &= 3 \neq 1,\\
\text{mcd}(6,10) &= 2 \neq 1,\\
\text{mcd}(15,10) &= 5 \neq 1.
\end{aligned}
\]
Por tanto, estamos ante el \textbf{segundo caso}.

\bigskip

\textbf{Paso 3: Sustitución con parámetro auxiliar \(t\)}

Tomamos el par \((6,15)\):
\[
6x + 15y = 3(2x + 5y) = 3t, \quad \text{donde } t = 2x + 5y.
\]

Reemplazamos en la ecuación original:
\[
3t + 10z = 7.
\]

\bigskip

\textbf{Paso 4: Resolución de la ecuación diofántica en \(t\) y \(z\)}

\textit{Divisiones:}
\[
10 / 3 = 3 \;\text{resto}\; \boxed{1}
\Rightarrow \text{mcd}(3,10) = \boxed{1}.
\]

\textbf{Tabla del Algoritmo Extendido:}
\[
\begin{array}{c|c|c|c|c}
i & r_i & q_i & u_i & v_i \\
\hline
0 & 3 & \text{-} & 0 & 1 \\
1 & \boxed{1} & 3 & 1 & -3 \\
\end{array}
\]

\[
\boxed{1} = 10 \cdot (1) + 3 \cdot (-3).
\]

\end{ejemplo}
\begin{ejemplo}[title={Ecuación diofántica con 3 incógnitas (Segundo caso)}]

Multiplicamos por \(7\):
\[
10 \cdot (7) + 3 \cdot (-21) = 7.
\]

Por tanto, una \textbf{solución particular} es:
\[
t_0 = -21, \quad z_0 = 7.
\]

\textbf{Solución general:}
\[
t = -21 + 10k, \quad z = 7 - 3k, \quad k \in \mathbb{Z}.
\]

\bigskip

\textbf{Paso 5: Sustitución en la ecuación definida en el Paso 3 para \(x\) e \(y\)}

Sabemos que:
\[
2x + 5y = t = -21 + 10k.
\]

\textit{Divisiones:}
\[
5 / 2 = 2 \;\text{resto}\; \boxed{1}
\Rightarrow \text{mcd}(2,5) = \boxed{1}.
\]

\textbf{Tabla del Algoritmo Extendido:}
\[
\begin{array}{c|c|c|c|c}
i & r_i & q_i & u_i & v_i \\
\hline
0 & 2 & \text{-} & 0 & 1 \\
1 & \boxed{1} & 2 & 1 & -2 \\
\end{array}
\]

\[
\boxed{1} = 2 \cdot (-2) + 5 \cdot (1).
\]

Multiplicamos por \((-21 + 10k)\):
\[
2 \cdot (-2)(-21 + 10k) + 5 \cdot (1)(-21 + 10k) = -21 + 10k.
\]

Por tanto, una \textbf{solución particular} es:
\[
x_0 = 42 - 20k, \quad y_0 = -21 + 10k.
\]

\bigskip

\textbf{Paso 6: Solución general de \(x, y, z\)}

\[
\begin{aligned}
x &= x_0 + 5k' = 42 - 20k + 5k',\\
y &= y_0 - 2k' = -21 + 10k - 2k',\\
z &= 7 - 3k,
\end{aligned}
\quad k, k' \in \mathbb{Z}.
\]

\bigskip

\textbf{Por tanto, las soluciones generales son:}
\[
\boxed{
\begin{aligned}
x &= 42 - 20k + 5k',\\
y &= -21 + 10k - 2k',\\
z &= 7 - 3k,
\end{aligned}
\quad k, k' \in \mathbb{Z}.
}
\]
\end{ejemplo}
%%%%%%%%%%%%%%%%%%%%%%%%%%%%%%%%%%%%%%%%%%%%%%%%%%%%%%%%%%%%

% =========================================================
% SECCIÓN: PROPIEDADES DE LA RELACIÓN DE CONGRUENCIA
% =========================================================

\section{La Relación de Congruencia en \(\mathbb{Z}\)}

\subsection{Propiedades}

\begin{procedimiento}[title={Definición de congruencia módulo \(m\)}]
Dados \(a,b \in \mathbb{Z}\) y \(m \in \mathbb{Z}\) con \(m \neq 0\), se dice que \(a\) es \textbf{congruente con} \(b\) \textbf{módulo \(m\)} si:
\[
a \equiv b \pmod{m} \;\;\Longleftrightarrow\;\; m \mid (a - b).
\]

De forma análoga, se dice que \(a\) y \(b\) \textbf{no son congruentes módulo \(m\)} si:
\[
a \not\equiv b \pmod{m} \;\;\Longleftrightarrow\;\; m \nmid (a - b).
\]
\end{procedimiento}

\begin{procedimiento}[title={Relación de equivalencia}]
La relación de congruencia módulo \(m\) es una \textbf{relación de equivalencia} en \(\mathbb{Z}\), ya que cumple:

\begin{enumerate}
    \item \textbf{Propiedad reflexiva:}  
    Para todo \(a \in \mathbb{Z}\),
    \[
    a \equiv a \pmod{m}.
    \]

    \item \textbf{Propiedad simétrica:}  
    Para cualesquiera \(a,b \in \mathbb{Z}\),
    \[
    a \equiv b \pmod{m} \;\;\Longleftrightarrow\;\; b \equiv a \pmod{m}.
    \]

    \item \textbf{Propiedad transitiva:}  
    Para cualesquiera \(a,b,c \in \mathbb{Z}\),
    \[
    a \equiv b \pmod{m} \text{ y } b \equiv c \pmod{m} \;\;\Longrightarrow\;\; a \equiv c \pmod{m}.
    \]
\end{enumerate}
\end{procedimiento}

\subsection{Clases de Equivalencia}

\begin{procedimiento}[title={Definición de clase de equivalencia módulo \(m\)}]
Sea \(a \in \mathbb{Z}\) y \(m \in \mathbb{Z}\setminus\{0\}\).  
La \textbf{clase de equivalencia de \(a\) módulo \(m\)} se define como:
\[
[a] = \{\, x \in \mathbb{Z} \mid x \equiv a \pmod{m} \,\}.
\]
\end{procedimiento}

\begin{procedimiento}[title={Propiedad de pertenencia}]
Todo entero pertenece a su clase:
\[
a \in [a].
\]
\end{procedimiento}

\begin{procedimiento}[title={Clases equivalentes}]
Para cualesquiera \(a,b \in \mathbb{Z}\):
\[
a \equiv b \pmod{m} \;\;\Longleftrightarrow\;\; [a] = [b].
\]
\[
a \not\equiv b \pmod{m} \;\;\Longleftrightarrow\;\; [a] \cap [b] = \varnothing.
\]
\end{procedimiento}

\begin{procedimiento}[title={Compatibilidad con suma y producto}]
Sean \(a,b,c,d \in \mathbb{Z}\) y \(m \in \mathbb{Z}\setminus\{0\}\).  
Si \(a \equiv c \pmod{m}\) y \(b \equiv d \pmod{m}\), entonces:
\[
a + b \equiv c + d \pmod{m}, \quad a \cdot b \equiv c \cdot d \pmod{m}.
\]

Esto significa que los elementos de una clase son intercambiables en operaciones módulo \(m\) sin alterar el resultado.
\end{procedimiento}

\begin{procedimiento}[title={Conjunto cociente \(\mathbb{Z}_m\)}]

El conjunto de todas las clases residuales módulo \(m\) se llama \textbf{conjunto cociente}:
\[
\mathbb{Z}_m = \{\, [0], [1], [2], \dots, [m-1] \,\}.
\]

En \(\mathbb{Z}_m\), las operaciones de suma y producto se definen mediante las clases:
\[
[a] + [b] = [a+b], \quad [a] \cdot [b] = [a \cdot b].
\]
\end{procedimiento}

\section{Inverso en Anillos Conmutativos}

\begin{procedimiento}[title={Inverso en \(\mathbb{Z}_m\)}]
Un elemento \([a] \in \mathbb{Z}_m\) es \textbf{invertible} (o \textbf{unidad}) si y solo si
\[
\text{mcd}(a,m) = 1.
\]

En tal caso, su inverso se obtiene mediante el Algoritmo Extendido de Euclides:
\[
[a]^{-1} = [u],
\]
donde \(u\) es el coeficiente de Bézout correspondiente a \(a\) en la combinación lineal
\[
a u + m v = 1.
\]

Cuidado, en la práctica \([u]\) estará en la columna de \(v_i\) porque corresponde a \(a\), que siempre será más pequeño que \(m\).
\end{procedimiento}

\begin{ejemplo}[title={Inverso de \([7]\) en \(\mathbb{Z}_{44}\)}]
\[
44 / 7 = 6 \;\text{resto } 2
\]
\[
7 / 2 = 3 \;\text{resto } \boxed{1}
\]

\[
\Rightarrow \text{mcd}(7,44) = 1 \;\;\text{y por tanto } [7] \text{ es unidad en } \mathbb{Z}_{44}.
\]

\textbf{Tabla del Algoritmo Extendido:}
\[
\begin{array}{c|c|c|c|c}
i & r_i & q_i & u_i & v_i \\
\hline
0 & 7 & - & 0 & 1 \\
1 & 2 & 6 & 1 & -6 \\
2 & \boxed{1} & 3 & -3 & 19 \\
\end{array}
\]

Por tanto:
\[
[7]^{-1} = [19] \in \mathbb{Z}_{44}.
\]
\end{ejemplo}

\begin{ejemplo}[title={\([15]\) en \(\mathbb{Z}_{171}\) no tiene inverso}]
\[
171 / 15 = 11 \;\text{resto } 6
\]
\[
15 / 6 = 2 \;\text{resto } 3
\]
\[
6 / 3 = 2 \;\text{resto } \boxed{0}
\]

\[
\text{mcd}(15,171) = 3 \neq 1 \;\;\Rightarrow [15] \text{ no es unidad en } \mathbb{Z}_{171}.
\]
\end{ejemplo}

\subsection{Función \(\varphi\) de Euler en \(\mathbb{Z}_m\)}

\begin{procedimiento}[title={Número de unidades en \(\mathbb{Z}_m\) y función \(\varphi(m)\)}]
Sea \(m \ge 2\). La función de Euler \(\varphi(m)\) representa el \textbf{número de unidades} en \(\mathbb{Z}_m\).

\bigskip
\textbf{Cálculo de \(\varphi(m)\):}
\begin{enumerate}
    \item Si \(p\) es primo y \(n \ge 1\):
    \[
    \varphi(p^n) = p^n - p^{\,n-1} = p^{n-1}(p-1).
    \]
    \item Si $a$ y $b$ son primos:
    \[
    \varphi(a \cdot b) = \varphi(a) \cdot \varphi(b).
    \]
\end{enumerate}
\end{procedimiento}

\begin{ejemplo}[title={Cálculo de \(\varphi(9000)\)}]
Factorizamos:
\[
9000 = 2^3 \cdot 3^2 \cdot 5^3
\]
\[
\varphi(9000) = \varphi(2^3) \cdot \varphi(3^2) \cdot \varphi(5^3)
\]

Calculamos \(\varphi\) de cada potencia de primo:

\[
\begin{aligned}
\varphi(2^3) &= 2^{3-1} (2-1) = 4\\
\varphi(3^2) &= 3^{2-1} (3-1) = 6\\
\varphi(5^3) &= 5^{3-1} (5-1) = 100
\end{aligned}
\]

\[
\varphi(9000) = 4 \cdot 6 \cdot 100 = 2400
\]

\textbf{Conclusión:} hay \(2400\) unidades en \(\mathbb{Z}_{9000}\).
\end{ejemplo}


\begin{procedimiento}[title={Cuándo \(\mathbb{Z}_m\) es un cuerpo}]
\(\mathbb{Z}_m\) es un \textbf{cuerpo} si y solo si \(m\) es primo.  
En ese caso, todos los elementos \([a] \neq [0]\) son unidades, y \(\varphi(m) = m-1\).
\end{procedimiento}

% =========================================================
% SECCIÓN: SISTEMAS DE ECUACIONES EN CONGRUENCIAS
% =========================================================
\section{Sistemas de Ecuaciones en Congruencias}


% =========================================================
% SUBSECCIÓN: CASO GENERAL
% =========================================================
\subsection{Ecuación de una congruencia lineal}

\begin{procedimiento}[title={Ecuación general}]
La ecuación
\[
a x \equiv b \pmod{m}
\]
tiene solución si y solo si
\[
\text{mcd}(a,m) \mid b.
\]

\begin{itemize}
    \item Si \(a = 1\), la solución general es:
    \[
    x = b + m k, \quad k \in \mathbb{Z}.
    \]
    \item En el caso general \(a \neq 1\):
    \[
    a x \equiv b \pmod{m},
    \]
    sea \(g = \text{mcd}(a,m)\), entonces existen coeficientes de Bézout \(u, v\) tales que:
    \[
    a u + m v = g.
    \]
    La solución general viene dada por:
    \[
    x = \frac{b}{g} \cdot u + \frac{m}{g} \cdot k, \quad k \in \mathbb{Z}.
    \]
\end{itemize}
\end{procedimiento}

% =========================================================
% EJEMPLO 1: ECUACIÓN INDIVIDUAL SIN MÚLTIPLES
% =========================================================
\begin{ejemplo}[title={Ejemplo 1: Ecuación \(102x \equiv 24 \pmod{318}\)}]

\textbf{Cálculo del MCD:}

\[
318 / 102 = 3 \;\text{resto}\; 12
\]
\[
102 / 12 = 8 \;\text{resto}\; \boxed{6}
\]
\[
12 / 6 = 2 \;\text{resto}\; 0
\]

\[
\text{mcd}(102, 318) = \boxed{6}
\]

Como \(6 \mid 24\), la ecuación tiene solución.

Dividimos toda la ecuación entre 6:
\[
17x \equiv 4 \pmod{53}
\]

\textbf{Cálculo del inverso de 17 módulo 53:}

\[
53 / 17 = 3 \;\text{resto}\; 2
\]
\[
17 / 2 = 8 \;\text{resto}\; \boxed{1}
\]

\[
\text{mcd}(17,53) = \boxed{1}
\]

\textbf{Tabla del Algoritmo Extendido de Euclides:}

\[
\begin{array}{c|c|c|c|c}
i & r_i & q_i & u_i & v_i \\
\hline
0 & 17 & - & 0 & 1 \\
1 & 2 & 3 & 1 & -3 \\
2 & \boxed{1} & 8 & -8 & 25 \\
\end{array}
\]

\[
\boxed{1} = 53 \cdot (-8) + 17 \cdot (25)
\]

Por tanto, el inverso de \([17]\) es \([25]\).

\[
x = \frac{4}{1} \cdot 25 + \frac{53}{1} \cdot k = 100 + 53k \equiv 47 + 53k, \quad k \in \mathbb{Z}.
\]
\end{ejemplo}

% =========================================================
% SUBSECCIÓN: SISTEMA DE CONGRUENCIAS
% =========================================================
\subsection{Resolución de un sistema de ecuaciones en congruencias}

\begin{procedimiento}[title={Método general}]
Para resolver un \textbf{sistema de ecuaciones en congruencias}:
\[
\begin{cases}
a_1 x \equiv b_1 \pmod{m_1} \\
a_2 x \equiv b_2 \pmod{m_2} \\
\vdots \\
a_n x \equiv b_n \pmod{m_n}
\end{cases}
\]

se sigue el siguiente procedimiento general:

\begin{enumerate}
    \item Se resuelve la primera congruencia individualmente, obteniendo una solución del tipo \(x = \text{(expresión con un parámetro entero)}\).
    \item Esa expresión se \textbf{sustituye} en la segunda congruencia, obteniendo una nueva ecuación que se resuelve del mismo modo.
    \item El proceso se \textbf{repite iterativamente}, sustituyendo el resultado anterior en la siguiente congruencia hasta llegar al final del sistema.
    \item Si en algún paso una de las ecuaciones no tiene solución, el sistema completo tampoco la tiene.
    \item Si todas las ecuaciones son compatibles, el resultado final será una única congruencia equivalente del tipo:
    \[
    x \equiv r \pmod{M},
    \]
    donde \(M\) es el módulo común que unifica todo el sistema.
\end{enumerate}

\end{procedimiento}

% =========================================================
% EJEMPLO 2: SISTEMA SIN SOLUCIÓN
% =========================================================
\begin{ejemplo}[title={Ejemplo 2: Sistema sin solución}]
\[
\begin{cases}
x \equiv 9 \pmod{12} \\
x \equiv 7 \pmod{20}
\end{cases}
\]

De la segunda ecuación:
\[
x = 7 + 20k, \quad k \in \mathbb{Z}.
\]

Sustituyendo en la primera:
\[
7 + 20k \equiv 9 \pmod{12}
\]
\[
20k \equiv 2 \pmod{12} \Rightarrow 8k \equiv 2 \pmod{12}
\]

\[
12 / 8 = 1 \;\text{resto}\; \boxed{4}, \quad 8 / 4 = 2 \;\text{resto}\; 0
\]

\[
\text{mcd}(8,12) = \boxed{4}
\]

Como \(4 \nmid 2\), el sistema \textbf{no tiene solución}.
\end{ejemplo}

% =========================================================
% EJEMPLO 3: SISTEMA CON SOLUCIÓN
% =========================================================
\begin{ejemplo}[title={Ejemplo 3: Sistema con solución}]
\[
\begin{cases}
x \equiv 9 \pmod{12} \\
x \equiv 5 \pmod{20}
\end{cases}
\]

De la segunda:
\[
x = 5 + 20k, \quad k \in \mathbb{Z}.
\]
Sustituyendo en la primera:
\[
5 + 20k \equiv 9 \pmod{12} \Rightarrow 20k \equiv 4 \pmod{12}
\]
\[
8k \equiv 4 \pmod{12}
\]

\[
12 / 8 = 1 \;\text{resto}\; \boxed{4}, \quad 8 / 4 = 2 \;\text{resto}\; 0
\]
\[
\text{mcd}(8,12) = \boxed{4}
\]
Como \(4 \mid 4\), la ecuación tiene solución.

Dividimos por 4:
\[
2k \equiv 1 \pmod{3}
\]

\textbf{Cálculo del inverso de 2 módulo 3:}

\[
3 / 2 = 1 \;\text{resto}\; \boxed{1}
\]
\[
\text{mcd}(2,3) = \boxed{1}
\]

\[
\begin{array}{c|c|c|c|c}
i & r_i & q_i & u_i & v_i \\
\hline
0 & 2 & - & 0 & 1 \\
1 & \boxed{1} & 1 & 1 & -1 \\
\end{array}
\]

\[
\boxed{1} = 3 \cdot (1) + 2 \cdot (-1)
\]

Por tanto, el inverso de \([2]\) es \([2]\) mismo (pues \(-1 \equiv 2 \pmod{3}\)).

\[
k = -1 + 3k', \quad k' \in \mathbb{Z}.
\]

Sustituyendo en \(x = 5 + 20k\):
\[
x = 5 + 20(-1 + 3k') = -15 + 60k' = 45 + 60k', \quad k' \in \mathbb{Z}.
\]

Por tanto, el sistema es equivalente a la ecuación:
\[
x \equiv 45 \pmod{60}.
\]
\end{ejemplo}

\section{Reducción de Potencias}

\begin{procedimiento}[title=Reducción de potencias]
Queremos calcular \(r_m(a^b)\), es decir, el resto de dividir \(a^b\) entre \(m\).

\begin{enumerate}
  \item \textbf{Reducir la base.}  
  Dividimos \(a\) entre \(m\) y el resto será la nueva base \(a'\).  
  Así obtenemos un nuevo \(r_m(a'^b)\).

  \item \textbf{Reducir el exponente.}  
  Calculamos \(\text{mcd}(a',m)\) y distinguimos dos casos:

  \begin{itemize}
    \item \textbf{Caso 1:} \(\text{mcd}(a',m)=1\)
    \begin{enumerate}
      \item[1.1)] Hallar \(l\) tal que \(a'^l \equiv 1 \pmod{m}\) calculando:
      \[
      a'^1 = a' \bmod m, \quad a'^2 = a'\cdot a' \bmod m, \quad a'^3 = a'^2 \cdot a' \bmod m, \ldots
      \]
      hasta que se cumpla \(a'^n \equiv 1 \pmod{m}\). Entonces \(l = n\). \\
      Para darnos pistas de cuánto vale $l$, podemos reducir el conjunto de valores que puede tomar $l$ si tenemos en cuenta que para este caso se cumple que $l|\varphi(m)$.
      \item[1.2)] Dividir \(b\) entre \(l\) y tomar el resto como nuevo exponente \(b'\).  
      Así obtenemos \(r_m(a'^{\,b'})\), que es mucho más sencillo de calcular.
    \end{enumerate}

    \item \textbf{Caso 2:} \(\text{mcd}(a',m)\neq1\)
    \begin{enumerate}
      \item[2.1)] Descomponer el módulo \(m\) en dos números: uno que tenga factores primos en común con la base (llamado \emph{primer número}) y otro que no (\emph{segundo número}).
      \item[2.2)] Llamamos \(a\) a la potencia actual.  
      Calculamos \(r_{\text{primer}}(a) = 0\), porque la base será múltiplo del primer número. Esto nos da la primera ecuación:
      \[
      x \equiv 0 \pmod{\text{primer número}}
      \]
      \item[2.3)] Calculamos \(r_{\text{segundo}}(a)\).  
      Como la base y el segundo número no comparten factores primos, su mcd es 1.  
      Reducimos la base y, si es necesario, el exponente usando el método de la \(l\), obteniendo así:
      \[
      x \equiv \text{resto} \pmod{\text{segundo número}}
      \]
      \item[2.4)] Formamos el sistema de congruencias:
      \[
      \begin{cases}
      x \equiv 0 \pmod{\text{primer número}}\\
      x \equiv \text{resto} \pmod{\text{segundo número}}
      \end{cases}
      \]
      y resolvemos. La solución menor que el módulo total se toma como el resto final (equivalente a escoger \(k'=0\)).
    \end{enumerate}
  \end{itemize}

\end{enumerate}
\end{procedimiento}

\begin{ejemplo}[title=Ejemplo 1: mcd es 1]
Cálculo de $r_{19}(45^{200})$

Reducimos la base:

\[
45/19 = 2\ \text{resto}\ 7
\]
\[
a' = 7 \quad\Rightarrow\quad r_{19}(45^{200}) = r_{19}(7^{200})
\]

Como \(\text{mcd}(7,19)=1\), buscamos un \(l\) tal que \(7^l = 1 \pmod{19}\):
\[
l \mid \varphi(19) \implies l \mid 18 \implies l \in \{1,2,3,6,9,18\}
\]

\begin{itemize}
  \item $7^1 = 7 \mod 19$
  \item $7^2 = 11 \mod 19$
  \item $7^3 = 77 \mod 19 \;\Rightarrow\; 7^3 = 1 \mod 19$
\end{itemize}

Por tanto, \(l = 3\).

Reducimos el exponente:

\[
200/3 = 66\ \text{resto}\ 2
\]
\[
b' = 2 \quad\Rightarrow\quad r_{19}(7^{200}) = r_{19}(7^2)
\]

\[
7^2 = 11 \pmod{19}
\]

\[
\boxed{r_{19}(45^{200}) = 11}
\]
\end{ejemplo}

\begin{ejemplo}[title=Ejemplo 2: mcd no es 1]
Calcular \(r_{36}(22^{12345})\).

La base ya está reducida. Comprobamos el mcd:

\[
\text{mcd}(22,36) = 2 \neq 1
\]

No existe \(l\) tal que \(22^l \equiv 1 \pmod{36}\).  
Descomponemos el módulo:

\[
36 = 4 \cdot 9
\]
donde 4 tiene primos en común con la base (22) y 9 no. Llamamos \(a = 22^{12345}\).

\begin{itemize}
    \item Es claro que $a$ es múltiplo de 4. Como \(r_4(a) = 0\).  
  Entonces $a$ satisface la primera ecuación del sistema:
  \[
  x \equiv 0 \pmod{4}
  \]
    \item Calculamos \(r_9(a) = r_9(4^{12345})\) porque \(22 \equiv 4 \pmod{9}\) y \(\text{mcd}(4,9) = 1\).  
  Hallamos \(l\) módulo 9:
  \[
    l \mid \varphi(9) \implies l \mid 6 \implies l \in \{1,2,3,6\}
    \]
  \begin{itemize}
    \item $4^1 = 4 \mod 9$
    \item $4^2 = 7 \mod 9$
    \item $4^3 = 28 \mod 9 \;\Rightarrow\; 4^3 = 1 \mod 9$
  \end{itemize}
  Por tanto \(l=3\).

- Reducimos el exponente:
\[
12345/3 = 4115 \text{ resto } 0
\]
\[
r_9(a) = r_9(4^{12345}) = r_9(4^0) = r_9(1) = 1
\]

Esto nos da la segunda ecuación:
\[
x \equiv 1 \pmod{9}
\]
\end{itemize}

Ahora resolvemos el sistema de congruencias:
\[
\begin{cases}
x \equiv 0 \pmod{4}\\
x \equiv 1 \pmod{9}
\end{cases}
\]

Escribimos \(x = 4k\), entonces:
\[
4k \equiv 1 \pmod{9}
\]

- Como \(\text{mcd}(4,9)=1\), la congruencia tiene solución.

\[
\begin{array}{c|c|c|c|c}
i & r_i & q_i & u_i & v_i \\
\hline
0 & 4 & - & 0 & 1 \\
1 & \boxed{1} & 2 & 1 & -2
\end{array}
\]

\[[1] = 4 \cdot (-2) + 9 \cdot 1\]

\end{ejemplo}

\begin{ejemplo}[title=Ejemplo 2: mcd no es 1]

Por tanto, el inverso de 4 módulo 9 es \(-2 \equiv 7 \pmod{9}\).

- Entonces \(k \equiv -2 \equiv 7 \pmod{9}\)  
\[
k = 7 + 9k',\quad k' \in \mathbb{Z}
\]

- y:
\[
x = 4k = 4(7 + 9k') = 28 + 36k',\quad k' \in \mathbb{Z}
\]

Finalmente, tomando \(k' = 0\):
\[
r_{36}(22^{12345}) = 28
\]
\end{ejemplo}

\section{Polinomios con Coeficientes en un Cuerpo}

\subsection{Divisibilidad de Polinomios}

\begin{procedimiento}[title=Divisibilidad de polinomios]
Sea \(\mathbb{K}\) un cuerpo y consideremos los polinomios \(f_1(x), f_2(x) \in \mathbb{K}[x]\).

Decimos que \(f_1(x)\) \textbf{divide} a \(f_2(x)\), y escribimos
\[
f_1(x) \mid f_2(x),
\]
si existe un polinomio \(g(x) \in \mathbb{K}[x]\) tal que
\[
f_2(x) = f_1(x) \cdot g(x).
\]
\end{procedimiento}

\subsection{Polinomios Asociados}

\begin{procedimiento}[title=Polinomios asociados]
Sea \(\mathbb{K}\) un cuerpo y consideremos los polinomios \(f(x), g(x) \in \mathbb{K}[x]\).

- Decimos que \(f(x)\) y \(g(x)\) están \textbf{asociados}, y escribimos
\[
f(x) \sim g(x),
\]
si uno \textbf{divide} al otro, es decir,
\[
f(x) \mid g(x) \quad \text{y} \quad g(x) \mid f(x).
\]

- Propiedad importante: Todo polinomio \(f(x) \neq 0\) tiene un \textbf{polinomio mónico asociado}.
\end{procedimiento}

\section{MCD y Coeficientes de Bézout de Polinomios}

\begin{procedimiento}[title=Cálculo del MCD y Coeficientes de Bézout]
Sea \(\mathbb{K}\) un cuerpo y consideremos dos polinomios \(f_1(x), f_2(x) \in \mathbb{K}[x]\).

El procedimiento para calcular el máximo común divisor (MCD) y los coeficientes de Bézout es el siguiente:

\begin{enumerate}
  \item Aplicamos el \textbf{Algoritmo Extendido de Euclides} para polinomios:
  \begin{itemize}
    \item Escogemos como \textbf{dividendo} el polinomio de mayor grado y como \textbf{divisor} el otro.
    \item Realizamos divisiones sucesivas hasta que el resto sea 0.
  \end{itemize}

  \item Durante el proceso, generamos una tabla con los siguientes datos:

\[
\begin{array}{c|c|c|c|c}
i & r_i(x) & q_i(x) & u_i(x) & v_i(x) \\
\hline
0 & f_2(x) & - & 0 & 1 \\
1 & r_1(x) & q_1(x) & 1 & -q_1(x) \\
2 & r_2(x) & q_2(x) & u_2(x) & v_2(x) \\
\vdots & \vdots & \vdots & \vdots & \vdots \\
n & 0 & q_n(x) & u_n(x) & v_n(x)
\end{array}
\]

  \item El \textbf{último resto no nulo} \(r_{n-1}(x)\) es un \textbf{máximo común divisor} de \(f_1(x)\) y \(f_2(x)\).  
  Para obtener el \textbf{mcd mónico}, multiplicamos por el inverso del coeficiente líder de \(r_{n-1}(x)\).

  \item Los polinomios \(u(x)\) y \(v(x)\) de Bézout se obtienen directamente de la columna correspondiente en la tabla, de manera que
  \[
  \text{mcd}(f_1,f_2) = u(x) \cdot f_1(x) + v(x) \cdot f_2(x).
  \]
\end{enumerate}
\end{procedimiento}

\begin{ejemplo}[title=Ejemplo de MCD y Coeficientes de Bézout en $\mathbb{Z}_{11}$]
Sea 
\[
f_1(x) = x^3 + 4x^2 + 4x + 5, \quad f_2(x) = 5x^2 + 2x + 5 \in \mathbb{Z}_{11}[x].
\]

Queremos calcular \(\text{mcd}(f_1(x),f_2(x))\).

Dividimos en \(\mathbb{Z}_{11}\):

\[
f_1(x)  / f_2(x) = -2x-5 \quad \text{resto } \boxed{2x-3}
\]

\[
f_2(x) / (2x-3) = 8x-2 \quad \text{resto } 0
\]

La tabla de Euclides extendido queda:

\[
\begin{array}{c|c|c|c|c}
i & r_i(x) & q_i(x) & u_i(x) & v_i(x) \\
\hline
0 & 5x^2+2x+5 & - & 0 & 1 \\
1 & \boxed{2x-3} & -2x-5 & 1 & 2x+5
\end{array}
\]

De la tabla, los coeficientes de Bézout son:

\[
f_1(x) \cdot (1) + f_2(x) \cdot (2x+5) = \boxed{2x-3}
\]

Para obtener el \textbf{MCD mónico}, multiplicamos por el inverso de 2 en \(\mathbb{Z}_{11}\), que es 6:

\[
f_1(x) \cdot (6) + f_2(x) \cdot (x+8) = x+4
\]

Por tanto:
\[
\text{mcd}(f_1(x), f_2(x)) = x+4
\]

Y 6 y \(x+8\) son unos \textbf{coeficientes de Bézout} para \(f_1(x)\) y \(f_2(x)\).
\end{ejemplo}


\section{Ecuaciones Diofánticas en \(\mathbb{K}[x]\)}

\begin{procedimiento}[title=Resolución de ecuaciones diofánticas en \(\mathbb{K}\)]
Sea \(\mathbb{K}\) un cuerpo y consideremos los polinomios \(a(x), b(x), c(x) \in \mathbb{K}[x]\).

Queremos resolver la ecuación diofántica:
\[
a(x) \cdot \alpha(x) + b(x) \cdot \beta(x) = c(x).
\]

El procedimiento es el siguiente:

\begin{enumerate}
  \item \textbf{Comprobar existencia de soluciones:}  
  La ecuación tiene solución si y solo si 
  \[
  \text{mcd}(a(x),b(x)) \mid c(x).
  \]

  \item \textbf{Calcular una solución particular:}  
  Se aplica el \textbf{Algoritmo Extendido de Euclides} para \(a(x)\) y \(b(x)\) y se obtiene una solución particular
  \((\alpha_0(x), \beta_0(x))\) que ya satisface
  \[
  a(x)\cdot\alpha_0(x) + b(x)\cdot\beta_0(x) = c(x).
  \]

  \item \textbf{Obtener la solución general:}  
  La solución general se expresa como
  \[
  \alpha(x) = \alpha_0(x) + \frac{b(x)}{g(x)} \cdot k(x), \quad
  \beta(x) = \beta_0(x) - \frac{a(x)}{g(x)} \cdot k(x),
  \]
  donde \(k(x) \in \mathbb{K}[x]\) es un polinomio arbitrario y \(g(x)\) es el mcd mónico de \(a(x)\) y \(b(x)\).
\end{enumerate}
\end{procedimiento}

\begin{ejemplo}[title=Resolución de una ecuación diofántica en $\mathbb{Z}_5$]
Sea
\[
a(x) = x^2+4x+3, \quad b(x) = x^2+2x+1, \quad c(x) = x^2+4 \in \mathbb{Z}_5[x].
\]

\textbf{Paso 1: Calcular el MCD de \(a(x)\) y \(b(x)\)}

\[
x^2+4x+3 / x^2+2x+1 = 1 \quad \text{resto } \boxed{2x+2}
\]

\[
x^2+2x+1 / 2x+2 = -2x-2 \quad \text{resto } 0
\]

\[
\text{mcd}(a(x),b(x)) = 2^{-1} \cdot \boxed{(2x+2)} = x+1
\]

\textbf{Paso 2: Comprobar divisibilidad por \(c(x)\)}

\[
x^2+4 / x+1 = x-1 \quad \text{resto } 0
\]

Por tanto \(x+1 \mid c(x)\), la ecuación tiene solución.

\textbf{Paso 3: Simplificar la ecuación}

\[
\frac{a(x)}{x+1} = x+3, \quad \frac{b(x)}{x+1} = x+1, \quad \frac{c(x)}{x+1} = x+4
\]

Ecuación simplificada:
\[
(x+3)\cdot\alpha(x) + (x+1)\cdot\beta(x) = x+4
\]

\textbf{Paso 4: Aplicar Algoritmo Extendido de Euclides}

\[
x+3 / x+1 = 1 \quad \text{resto } \boxed{2}
\]

Tabla de Euclides extendido:

\[
\begin{array}{c|c|c|c|c}
i & r_i(x) & q_i(x) & u_i(x) & v_i(x) \\
\hline
0 & x+1 & - & 0 & 1 \\
1 & \boxed{2} & 1 & 1 & -1
\end{array}
\]

\[
(x+3)\cdot (1) + (x+1)\cdot(-1) = \boxed{2}
\]

\[
\text{mcd}(x+3,x+1) = 2^{-1} \cdot \boxed{2} = 1
\]

Multiplicamos toda la ecuación por \(2^{-1} = 3\) en \(\mathbb{Z}_5\):

\[
(x+3)\cdot (3) + (x+1)\cdot(-3) = 1
\]

Multiplicamos por \(x+4\) para ajustar a la ecuación original:

\[
(x+3)\cdot (3x-3) + (x+1)\cdot (-3x+3) = x+4
\]

\textbf{Paso 5: Solución particular}

\[
\alpha_0(x) = 3x-3 = 3x+2, \quad \beta_0(x) = -3x+3 = 2x+3
\]

\end{ejemplo}

\begin{ejemplo}[title=Resolución de una ecuación diofántica en $\mathbb{Z}_5$]

\textbf{Paso 6: Solución general}

\[
\alpha(x) = (3x+2) + (x+1) k(x), \quad
\beta(x) = (2x+3) - (x+3) k(x), \quad k(x) \in \mathbb{Z}_5[x]
\]

Así, todas las soluciones de la ecuación diofántica están dadas por:

\[
\alpha(x) = 3x+2+(x+1)k(x), \quad
\beta(x) = 2x+3+4(x+3)k(x), \quad k(x) \in \mathbb{Z}_5[x].
\]
\end{ejemplo}

\section{MCM de Polinomios}

\begin{procedimiento}[title=Mínimo común múltiplo de polinomios]
Sea \(\mathbb{K}\) un cuerpo y consideremos polinomios \(f_1(x), f_2(x), \dots, f_n(x) \in \mathbb{K}[x]\).

- El \textbf{mínimo común múltiplo} (MCM) de \(f_1(x), f_2(x), \dots, f_n(x)\) es el polinomio mónico de \(\mathbb{K}[x]\) de menor grado que es divisible por todos los polinomios \(f_i(x)\).

\bigskip

- Propiedad importante (para dos polinomios):
\[
f_1(x) \cdot f_2(x) \sim \text{mcd}(f_1(x),f_2(x)) \cdot \text{mcm}(f_1(x),f_2(x))
\]
\end{procedimiento}

\section{Raíces y Multiplicidad de Polinomios}

\begin{procedimiento}[title=Raíces y multiplicidad]
Sea \(f(x) \in \mathbb{K}[x]\), donde \(\mathbb{K}\) es un cuerpo.

\begin{enumerate}
  \item \textbf{Teorema del resto:}  
  Para dividir \(f(x)\) entre \((x-\alpha)\):
  \[
  f(x) = q(x)(x-\alpha) + f(\alpha)
  \]
  y podemos usar \textbf{Ruffini} para realizar la división.

  \item \textbf{Raíz:}  
  Un elemento \(\alpha \in \mathbb{K}' \subseteq \mathbb{K}\) es raíz de \(f(x)\) si
  \[
  f(\alpha) = 0 \quad \Longleftrightarrow \quad (x-\alpha) \mid f(x).
  \]

  \item \textbf{Multiplicidad:}  
  La multiplicidad de una raíz \(\alpha\) es el mayor entero \(m \ge 1\) tal que
  \[
  (x-\alpha)^m \mid f(x).
  \]

  \item \textbf{Cálculo práctico:}  
  \begin{enumerate}
    \item Evaluamos \(f(\alpha)\) para encontrar una raíz \(\alpha\).
    \item Aplicamos Ruffini sucesivas veces con \(x-\alpha\) hasta que el resto deje de ser cero.  
          El número de veces que podemos dividir \(f(x)\) entre \(x-\alpha\) sin resto es la multiplicidad \(m\).
    \item El polinomio resultante \(g(x)\) (ya sin la raíz \(\alpha\) repetida) se vuelve a evaluar para encontrar nuevas raíces y sus multiplicidades.
    \item Repetimos el proceso hasta encontrar todas las raíces en \(\mathbb{K}[x]\).
  \end{enumerate}
\end{enumerate}
\end{procedimiento}

\begin{ejemplo}[title=Raíces y multiplicidades en $\mathbb{Z}_3$]
Sea
\[
f(x) = x^4 + x^3 + 2x^2 + x + 1 \in \mathbb{Z}_3[x].
\]

\textbf{Paso 1: Evaluar posibles raíces}

\[
f(0) = 1 \neq 0 \quad\Rightarrow\quad 0 \text{ no es raíz}
\]

\[
f(1) = 1+1+2+1+1 = 6 \equiv 0 \pmod{3} \quad\Rightarrow\quad \alpha=1 \text{ es raíz}
\]

\textbf{Paso 2: Comprobar multiplicidad con Ruffini}

\[
\begin{array}{r|rrrrr}
  & 1 & 1 & 2 & 1 & 1 \\
1 &   & 1 & 2 & 1 & 2 \\
\hline
  & 1 & 2 & 1 & 2 & \boxed{0} \\
1 &   & 1 & 0 & 1 \\
\hline
  & 1 & 0 & 1  & \boxed{0} \\
1 &   & 1 & 1 \\
\hline
  & 1 & 1 & 2
\end{array}
\]

No se puede seguir dividiendo sin resto. Por tanto, la multiplicidad de \(\alpha=1\) es \(\boxed{2}\).

\textbf{Paso 3: Polinomio resultante}

\[
f(x) = (x-1)^{\boxed{2}} \cdot (x^2 + 1)
\]

\textbf{Paso 4: Comprobar raíces de \(g(x) = x^2+1\) en $\mathbb{Z}_3$}

\[
g(0) = 1 \neq 0, \quad g(1) = 2 \neq 0, \quad g(2) = 2 \neq 0
\]

No hay más raíces en \(\mathbb{Z}_3\).

\textbf{Conclusión:}  
El polinomio \(f(x)\) tiene una raíz \(\alpha = 1\) con multiplicidad \(2\).
\end{ejemplo}

\section{Polinomios Irreducibles}

\begin{procedimiento}[title=Polinomios irreducibles]
Sea \(\mathbb{K}\) un cuerpo y \(f(x) \in \mathbb{K}[x]\).  

- Un polinomio \(f(x)\) es \textbf{irreducible} si no puede factorizarse como producto de polinomios de grado menor en \(\mathbb{K}[x]\).  

\medskip
\textbf{Criterio según el grado:}
\begin{itemize}
  \item \textbf{Grado 1:} Todos los polinomios son irreducibles.
  \item \textbf{Grado 2 o 3:} Irreducible si no tiene raíces en \(\mathbb{K}\).
  \item \textbf{Grado 4 o más:} No siempre irreducible si no tiene raíces en \(\mathbb{K}\).
\end{itemize}

\medskip
\textbf{Ejemplos de polinomios mónicos irreducibles:}

\vspace{2mm}
\begin{tabular}{c|l}
\textbf{Grado} & \textbf{Polinomios mónicos irreducibles en \(\mathbb{Z}_2[x]\)} \\
\hline
1 & $x$, $x+1$ \\
2 & $x^2+x+1$ \\
3 & $x^3+x+1$, $x^3+x^2+1$
\end{tabular}

\vspace{2mm}

\begin{tabular}{c|l}
\textbf{Grado} & \textbf{Polinomios mónicos irreducibles en \(\mathbb{Z}_3[x]\)} \\
\hline
1 & $x$, $x+1$, $x+2$ \\
2 & $x^2+1$, $x^2+x+2$, $x^2+2x+2$ \\
\end{tabular}
\end{procedimiento}

\begin{procedimiento}[title=Procedimiento para factorizar un polinomio en \(\mathbb{K}\)]
Sea \(f(x) \in \mathbb{K}[x]\). Para factorizarlo en producto de polinomios irreducibles:

\begin{enumerate}
  \item \textbf{Según el grado:}
  \begin{itemize}
    \item Si \(\deg(f) = 1\): es irreducible.
    \item Si \(\deg(f) = 2\) o \(3\): es irreducible si no tiene raíces en \(\mathbb{K}\).
    \item Si \(\deg(f) \ge 4\): continuar con el siguiente paso.
  \end{itemize}

  \item \textbf{Comprobación de raíces:}  
  Evaluar \(f(a)\) para todos los \(a \in \mathbb{K}\).  
  Cada raíz \(a\) encontrada da lugar a un factor \((x - a)\).  
  Si no tiene raíces, no tiene factores de grado 1 ni de grado \(n-1\).

  \item \textbf{Comprobación de factores de mayor grado:}  
  Si no hay raíces, probar divisiones de \(f(x)\) entre los \textbf{polinomios mónicos irreducibles conocidos} (de las tablas) de grado 2, luego 3, etc.  
  Si alguno divide exactamente, obtener el cociente y repetir el proceso con él.

  \item \textbf{Finalización:}  
  Continuar dividiendo hasta que todos los factores obtenidos sean irreducibles en \(\mathbb{K}[x]\).  
  El resultado final será la factorización completa:
  \[
  f(x) = p_1(x) \, p_2(x) \, \cdots \, p_k(x)
  \]
  donde cada \(p_i(x)\) es irreducible.
\end{enumerate}
\end{procedimiento}


\begin{ejemplo}[title=Ejemplo: Polinomio Irreducible]

Comprobar si \( f(x) = x^6 + x^3 + 1 \) es irreducible en \(\mathbb{Z}_2[x]\).

\medskip
\textbf{1. Raíces (factores de grado 1):}

Probamos si \(f(x)\) tiene raíces en \(\mathbb{Z}_2\):

\[
\begin{array}{lcl}
f(0) &=& 1 \neq 0 \\
f(1) &=& 1 + 1 + 1 = 3 \equiv 1 \neq 0
\end{array}
\]

Por tanto, \(f(x)\) no tiene raíces en \(\mathbb{Z}_2\)  
\(\Rightarrow\) \(f(x)\) no tiene factores de grado 1 ni de grado 5.

\medskip
\textbf{2. Factores de grado 2:}

De la tabla de polinomios mónicos irreducibles en \(\mathbb{Z}_2[x]\), sólo tenemos:
\[
x^2 + x + 1
\]

Dividimos:
\[
x^6 + x^3 + 1 \; / \; (x^2 + x + 1) = x^4 + x^3 \quad \text{resto } \boxed{1}
\]

\(\Rightarrow\) No divide, luego \(f(x)\) no tiene factores de grado 2 ni de grado 4.

\medskip
\textbf{3. Factores de grado 3:}

Polinomios mónicos irreducibles de grado 3 en \(\mathbb{Z}_2[x]\):
\[
x^3 + x + 1, \quad x^3 + x^2 + 1
\]

\[
\begin{array}{lcl}
x^6 + x^3 + 1 \; / \; (x^3 + x + 1) &=& x^3 + x \quad \text{resto } x^2 + x + 1 \\
x^6 + x^3 + 1 \; / \; (x^3 + x^2 + 1) &=& x^3 + x^2 + x + 1 \quad \text{resto } x
\end{array}
\]

En ambos casos el resto es distinto de 0, por tanto ninguno divide.

\medskip
\textbf{Conclusión:}
\[
f(x) \text{ no tiene factores de grado } 1, 2 \text{ ni } 3 \Rightarrow \boxed{f(x) \text{ es irreducible en } \mathbb{Z}_2[x].}
\]
\end{ejemplo}

\begin{ejemplo}[title=Ejemplo: Polinomio Reducible]

Factorizar \( f(x) = x^5 + x^3 + x^2 + x + 1 \) en \(\mathbb{Z}_3[x]\) como producto de polinomios irreducibles.

\medskip
\textbf{1. Raíces (factores de grado 1):}

Comprobamos si \(f(x)\) tiene raíces en \(\mathbb{Z}_3\):

\[
\begin{array}{lcl}
f(0) &=& 1 \neq 0 \\
f(1) &=& 1 + 1 + 1 + 1 + 1 = 5 \equiv 2 \neq 0 \\
f(2) &=& 2^5 + 2^3 + 2^2 + 2 + 1 = 32 + 8 + 4 + 2 + 1 = 47 \equiv 2 \neq 0
\end{array}
\]

Por tanto, \(f(x)\) no tiene raíces en \(\mathbb{Z}_3\)  
\(\Rightarrow\) no tiene factores de grado 1 ni de grado 4.

\medskip
\textbf{2. Factores de grado 2:}

Polinomios mónicos irreducibles de grado 2 en \(\mathbb{Z}_3[x]\):
\[
x^2 + 1,\quad x^2 + x + 2,\quad x^2 + 2x + 2
\]

Probamos dividir \(f(x)\) entre cada uno:

\[
\begin{array}{lcl}
f(x) / (x^2 + 1) &=& x^3 + 1 \quad \text{resto } \boxed{x} \quad \Rightarrow \text{no divide} \\[4pt]
f(x) / (x^2 + x + 2) &=& x^3 + 2x^2 \quad \text{resto } \boxed{x + 1} \quad \Rightarrow \text{no divide} \\[4pt]
f(x) / (x^2 + 2x + 2) &=& x^3 + x^2 + 2 \quad \text{resto } \boxed{0} \quad \Rightarrow \text{sí divide}
\end{array}
\]

Por tanto:
\[
f(x) = (x^2 + 2x + 2)(x^3 + x^2 + 2)
\]

\medskip
\textbf{3. Comprobamos si $g(x) = x^3 + x^2 + 2$ es irreducible:}

\[
\begin{array}{lcl}
g(0) &=& 2 \neq 0 \\
g(1) &=& 1 + 1 + 2 = 4 \equiv 1 \neq 0 \\
g(2) &=& 8 + 4 + 2 = 14 \equiv 2 \neq 0
\end{array}
\]

\(\Rightarrow g(x)\) no tiene raíces en \(\mathbb{Z}_3\), por tanto es \textbf{irreducible}.

\medskip
\textbf{Conclusión:}
\[
\boxed{f(x) = (x^2 + 2x + 2)(x^3 + x^2 + 2)}
\]
donde ambos factores son irreducibles en \(\mathbb{Z}_3[x]\).
\end{ejemplo}

\begin{procedimiento}[title=Caso particular de factorización con exponente común en $\mathbb{Z}_p\text{[$x$]}$]
Sea
\[
f(x)=a_n\,x^{n\cdot p}+a_{n-1}\,x^{(n-1)\cdot p}+\cdots + a_1\,x^{p}+a_0 \in \mathbb{Z}_p[x].
\]
\[
\text{Entonces }
f(x) = 
(a_n\,x^{n}+a_{n-1}\,x^{n-1}+\cdots+a_1\,x+a_0)^p.
\]
\end{procedimiento}

\begin{ejemplo}[title=Ejemplo de reducción de exponentes]
Sea
\[
f(x)=x^{42}+2x^{28}+3x^{21}+x^{14}+3x^{7}+2 \in \mathbb{Z}_7[x].
\]
\[
\text{Entonces }
f(x) = (
x^{6}+2x^{4}+3x^{3}+x^{2}+3x+2)^7.
\]
\end{ejemplo}

\section{Relación de Congruencia en \(\mathbb{K}[x]\)}

\subsection{Propiedades}

\begin{procedimiento}[title=Definición de congruencia módulo \(m(x)\)]
Dados \(f(x), g(x) \in \mathbb{K}[x]\), se dice que \(f(x)\) es \textbf{congruente} con \(g(x)\) módulo \(m(x)\) si:
\[
f(x) \equiv g(x) \pmod{m(x)} \;\;\Longleftrightarrow\;\; m(x) \mid \big(f(x) - g(x)\big)
\]

De forma análoga, se dice que \(f(x)\) y \(g(x)\) \textbf{no son congruentes} módulo \(m(x)\) si:
\[
f(x) \not\equiv g(x) \pmod{m(x)} \;\;\Longleftrightarrow\;\; m(x) \nmid \big(f(x) - g(x)\big)
\]

Si \(r(x)\) es el resto de dividir \(f(x)\) entre \(m(x)\), entonces:
\[
f(x) \equiv r(x) \pmod{m(x)}
\]
\end{procedimiento}

\begin{procedimiento}[title=Relación de equivalencia]
La relación de congruencia módulo \(m(x)\) es una \textbf{relación de equivalencia} en \(\mathbb{K}[x]\), ya que cumple las siguientes propiedades:

\begin{itemize}
    \item \textbf{Reflexiva:} \(f(x) \equiv f(x) \pmod{m(x)}\)
    \item \textbf{Simétrica:} si \(f(x) \equiv g(x) \pmod{m(x)}\), entonces \(g(x) \equiv f(x) \pmod{m(x)}\)
    \item \textbf{Transitiva:} si \(f(x) \equiv g(x) \pmod{m(x)}\) y \(g(x) \equiv h(x) \pmod{m(x)}\), entonces \(f(x) \equiv h(x) \pmod{m(x)}\)
\end{itemize}
\end{procedimiento}

\subsection{Clases de Equivalencia}

\begin{procedimiento}[title=Definición de clase de equivalencia módulo \(m(x)\)]
La \textbf{clase de equivalencia} de un polinomio \(f(x)\) módulo \(m(x)\) es el conjunto:
\[
[f(x)] = \{\, g(x) \in \mathbb{K}[x] \;|\; f(x) \equiv g(x) \pmod{m(x)} \,\}
\]
\end{procedimiento}

\begin{procedimiento}[title=Conjunto cociente ]
El \textbf{conjunto cociente} es el conjunto formado por todas las clases de equivalencia módulo \(m(x)\), y se denota por:
\[
\mathbb{K}[x]_{m(x)} = \{\, [f(x)] \;|\; f(x) \in \mathbb{K}[x] \,\}
\]

En \(\mathbb{K}[x]_{m(x)}\), las operaciones de suma y producto se definen mediante las clases:
\[
[f(x)] + [g(x)] = [\,f(x) + g(x)\,], \qquad [f(x)] \cdot [g(x)] = [\,f(x) \cdot g(x)\,].
\]
\end{procedimiento}


\begin{ejemplo}
\[
\mathbb{Z}_2[x]_{x^3+x+1} = \{[0],[1],[x],[x+1],[x^2],[x^2+1],[x^2+x],[x^2+x+1]\}
\]
\end{ejemplo}

\begin{procedimiento}[title=Cardinal del conjunto cociente]
El número de elementos del conjunto cociente \(\mathbb{Z}_p[x]_{m(x)}\) viene dado por:
\[
\bigl|\mathbb{Z}_p[x]_{m(x)}\bigr| = p^{\operatorname{grado}(m(x))}.
\]
\end{procedimiento}

\begin{ejemplo}
Sea \(A = \mathbb{Z}_5[x]_{x^3+2x+1}\).  
El número de elementos de \(A\) es:
\[
|A| = 5^3 = 125.
\]
Como \(x^3+2x+1\) es irreducible en \(\mathbb{Z}_5[x]\), \(A\) es un \textbf{cuerpo} de 125 elementos.
\end{ejemplo}

\subsubsection{Inverso de un polinomio en \(\mathbb{K}[x]_{m(x)}\)}

\begin{procedimiento}[title=Inverso en \(\mathbb{K}\)]
Un elemento \([f(x)] \in \mathbb{K}[x]_{m(x)}\) es \textbf{invertible} (o \textbf{unidad}) si y sólo si
\[
\operatorname{mcd}(f(x), m(x)) = 1.
\]

En tal caso, su inverso se obtiene mediante el \textbf{Algoritmo Extendido de Euclides}:
\[
[f(x)]^{-1} = [u(x)],
\]
donde \(u(x)\) es el coeficiente de Bézout correspondiente a \(f(x)\) en la combinación lineal
\[
f(x)\cdot u(x) + m(x)\cdot v(x) = 1.
\]

\textbf{Cuidado:} en la práctica, \([u(x)]\) se encuentra en la columna de \(v_i(x)\), ya que corresponde al polinomio \(f(x)\), que siempre será de menor grado que \(m(x)\).
\end{procedimiento}


\begin{ejemplo}
Sea \(A = \mathbb{Z}_5[x]_{x^3+2x+1}\) y \(f(x) = 3x^2 + x + 4\).  
¿Tiene \(f(x)\) inverso en \(A\)?

\[
x^3 + 2x + 1 \;/\; (3x^2 + x + 4) = 2x + 1 \quad \text{resto } 3x + 2
\]
\[
3x^2 + x + 4 \;/\; (3x + 2) = x + 3 \quad \text{resto } \boxed{3}
\]

En polinomios, basta parar de dividir cuando el resto es 0 o una constante.  
Por lo tanto, \(3\) es un mcd de \(f(x)\) y \(m(x)\).

Para obtener el mcd mónico, multiplicamos dicho mcd por el inverso del coeficiente líder de \(f(x)\).  
Como \(3^{-1} = 2\) en \(\mathbb{Z}_5\):
\[
3 \cdot 3^{-1} = 3 \cdot 2 = 1 = \operatorname{mcd}(f(x), m(x))
\]
Por lo tanto, \(f(x)\) tiene inverso en \(A\).

\vspace{0.5em}
\noindent\textbf{Algoritmo de Euclides Extendido:}
\[
\begin{array}{c|c|c|c|c}
i & r_i(x) & q_i(x) & u_i(x) & v_i(x) \\ \hline
0 & 3x^2 + x + 4 & - & 1 & 0 \\[4pt]
1 & 3x + 2 & 2x + 1 & 1 & -2x-1 \\[4pt]
2 & \boxed{3} & x + 3 & -x - 3 & 2x^2 + 2x + 4
\end{array}
\]

De la última fila se obtiene:
\[
f(x)\,(2x^2 + 2x + 4) + m(x)\,(-x + 3) = \boxed{3}.
\]

Multiplicamos por \(3^{-1} = 2\) en \(\mathbb{Z}_5\):
\[
f(x)\,(4x^2 + 4x + 3) + m(x)\,(-2x - 1) = 1 = \operatorname{mcd}(f(x), m(x)).
\]

Por tanto:
\[
[f(x)]^{-1} = [4x^2 + 4x + 3].
\]
\end{ejemplo}

\subsubsection{Cuerpos finitos mediante conjuntos cociente}

\begin{procedimiento}[title=Determinación de cuándo \(\mathbb{K}\) es un cuerpo]
El conjunto cociente \(\mathbb{K}[x]_{m(x)}\) es un \textbf{cuerpo} si y sólo si el polinomio \(m(x)\) es \textbf{irreducible} en \(\mathbb{K}[x]\).

\end{procedimiento}

\begin{ejemplo}[title=Ejemplo: Comprobación de que \(A\) es un cuerpo]
Sea \(A = \mathbb{Z}_5[x]_{x^3+2x+1}\). ¿Es un cuerpo?

\vspace{0.5em}

Comprobamos si \(m(x) = x^3 + 2x + 1\) es irreducible en \(\mathbb{Z}_5[x]\) verificando raíces en \(\mathbb{Z}_5\):

\[
\begin{aligned}
m(0) &= 1 \neq 0, \\
m(1) &= 1 + 2 + 1 = 4 \neq 0, \\
m(2) &= 8 + 4 + 1 = 3 \neq 0, \\
m(3) &= 27 + 6 + 1 = 4 \neq 0, \\
m(4) &= 64 + 8 + 1 = 3 \neq 0.
\end{aligned}
\]

Como \(m(x)\) no tiene raíces en \(\mathbb{Z}_5\), no tiene factores de grado 1, ni de grado 2.  
Por lo tanto, \(m(x)\) es irreducible y \(A\) es un \textbf{cuerpo}.
\end{ejemplo}

\begin{procedimiento}[title=Cómo construir un cuerpo con \(m\) elementos]

Para construir un cuerpo finito con \(m\) elementos, es necesario que \(m\) sea una \textbf{potencia de un primo}, ya que si no, no existe cuerpo con \(m\) elementos:

\begin{enumerate}
    \item Factorizamos \(m\) como potencia de un primo: \(m = p^n\).  
    \item Buscamos un polinomio \(m(x)\) irreducible de grado \(n\) en \(\mathbb{Z}_p[x]\).  
\end{enumerate}

El conjunto cociente
\[
\mathbb{Z}_p[x]_{m(x)}
\]
resultante será un \textbf{cuerpo} de \(m = p^n\) elementos.

\end{procedimiento}

\begin{ejemplo}[title=Ejemplo: Construcción de un cuerpo finito]
Se pide construir un cuerpo con \(16\) elementos.

\begin{enumerate}
    \item Factorizamos \(16 = 2^4\). Tomamos \(p = 2\) y \(n = 4\).  
    \item Buscamos un polinomio irreducible de grado \(4\) en \(\mathbb{Z}_2[x]\). Un ejemplo es:
    \[
    m(x) = x^4 + x + 1.
    \]
    Comprobamos que no tiene raíces en \(\mathbb{Z}_2\):
    \[
    m(0) = 1 \neq 0, \quad m(1) = 1 + 1 + 1 = 1 \neq 0.
    \]
    Como no tiene raíces (no tiene factores de grado \(1\)), tampoco tiene factores de grado \(3\).
    
    Entonces comprobamos que no es divisible entre ningún polinomio mónico irreducible de grado \(2\) en \(\mathbb{Z}_2\) y, efectivamente, \(m(x)\) es irreducible.
\end{enumerate}

Por lo tanto, el conjunto cociente
\(
\mathbb{Z}_2[x]_{x^4+x+1}
\)
es un \textbf{cuerpo} de \(16\) elementos.
\end{ejemplo}

\section{Ecuaciones en congruencia en \(\mathbb{K}[x]\)}

\begin{procedimiento}[title=Resolución de ecuaciones en congruencia]
Las ecuaciones en congruencia en \(\mathbb{K}[x]\) se resuelven de manera análoga a \(\mathbb{Z}\).  
Dada una ecuación
\[
a(x) f(x) \equiv b(x) \pmod{m(x)},
\]
primero calculamos \(\operatorname{mcd}(a(x), m(x))\):
\begin{itemize}
    \item Si \(\operatorname{mcd}(a(x), m(x)) \nmid b(x)\), la ecuación no tiene solución.
    \item Si \(\operatorname{mcd}(a(x), m(x)) \mid b(x)\), podemos simplificar dividiendo entre el mcd y resolver la ecuación reducida.
\end{itemize}
\end{procedimiento}

\begin{ejemplo}[title=Ejemplo 1: Sin solución]
En \(\mathbb{Z}_2[x]\) resolvemos:
\[
x f(x) \equiv x+1 \pmod{x^2+x}.
\]

Calculamos el mcd:
\[
\operatorname{mcd}(x, x^2+x) = x.
\]

Comprobamos si divide al término independiente:
\[
x \nmid (x+1) \quad \Rightarrow \quad \text{no tiene solución.}
\]
\end{ejemplo}

\begin{ejemplo}[title=Ejemplo 2: Con solución]
En \(\mathbb{Z}_2[x]\) resolvemos:
\[
(x+1) f(x) \equiv x^2+1 \pmod{x^2+x}.
\]

Calculamos el mcd:
\[
\operatorname{mcd}(x+1, x^2+x) = x+1.
\]

Comprobamos divisibilidad:
\[
x+1 \mid x^2+1 \quad \Rightarrow \quad \text{sí tiene solución.}
\]

Dividimos la ecuación entre \(x+1\):
\[
f(x) \equiv \frac{x^2+1}{x+1} \equiv x+1 \pmod{x} \quad \Rightarrow \quad f(x) \equiv 1 \pmod{x}.
\]

Por tanto, la solución general es:
\[
f(x) = 1 + x k(x), \quad k(x) \in \mathbb{Z}_2[x].
\]
\end{ejemplo}

\section{Ecuaciones de grado 1 en \(\mathbb{K}[x]_{m(x)}\)}

\begin{procedimiento}[title=Resolución de ecuaciones de grado 1]
Si tenemos una ecuación de la forma
\[
a(x)f(x) = b(x),
\]
el procedimiento es el siguiente:

\begin{enumerate}
    \item Comprobar si \(a(x)\) tiene inverso en \(\mathbb{K}[x]_{m(x)}\).  
    \begin{itemize}
        \item Si \(\operatorname{mcd}(a(x), m(x)) = 1\), multiplicamos por el inverso de \(a(x)\) y resolvemos normalmente.
        \item Si \(\operatorname{mcd}(a(x), m(x)) \neq 1\), transformamos la ecuación a congruencia
        \[
        a(x) f(x) \equiv b(x) \pmod{m(x)}
        \]
        y verificamos si \(\operatorname{mcd}(a(x), m(x))\) divide a \(b(x)\). Si no, no hay solución; si sí, dividimos entre el mcd y continuamos.
    \end{itemize}
    \item Hallar la solución general \(f(x)\) como congruencia.  
    \item Ajustar obligatoriamente la solución para que \(\deg(f(x)) < \deg(m(x))\).
    
    Obtenemos todos los valores de \(k(x)\) que hagan que la expresión en la que aparece tenga el mismo grado que \(f(x)\) y obtenemos cada solución de \(f(x)\) para cada valor de \(k(x)\).
\end{enumerate}
\end{procedimiento}

\begin{ejemplo}[title=Ejemplo: Resolución de ecuación de grado 1]
En \(\mathbb{Z}_2[x]_{x^3+x^2}\) resolver:
\[
x^2 f(x) \equiv x^2.
\]

\textbf{Paso 1:} Comprobamos si \(x^2\) tiene inverso:
\[
\operatorname{mcd}(x^2, x^3+x^2) = x^2 \neq 1 \quad \Rightarrow \text{no tiene inverso.}
\]

\textbf{Paso 2:} Transformamos a congruencia:
\[
x^2 f(x) \equiv x^2 \pmod{x^3+x^2}.
\]

\(\operatorname{mcd}(x^2, x^3+x^2) = x^2\) y \(x^2 \mid x^2 \) → tiene solución. Dividimos entre \(x^2\):
\[
f(x) \equiv 1 \pmod{x+1} \quad \Rightarrow \quad f(x) = 1 + (x+1) k(x), \quad k(x) \in \mathbb{Z}_2[x].
\]

\textbf{Paso 3:} Buscamos los \(f(x)\) cuyo grado es menor que el grado de \(m(x) = x^3+x^2\), es decir, \(\deg(f(x)) \le 2\).  
Esto implica que \(\deg(k(x)) \le 1\), entonces los posibles valores son:
\[
k(x) = 0, 1, x, x+1.
\]

\textbf{Soluciones:}
\[
\begin{aligned}
k(x) = 0 &\implies f(x) = 1, \\
k(x) = 1 &\implies f(x) = x, \\
k(x) = x &\implies f(x) = x^2 + x + 1, \\
k(x) = x+1 &\implies f(x) = x^2.
\end{aligned}
\]
\end{ejemplo}


\end{document}
